% =============================================================================
% Préambule — Mémoire BTP RAG (rapport_btp)
% Palette : bleu clair institutionnel
% =============================================================================

% -----------------------------------------------------------------------------
% Encodage, langue, typographie
% -----------------------------------------------------------------------------
\usepackage[T1]{fontenc}
\usepackage[utf8]{inputenc}
\usepackage[english, main=french]{babel}
\usepackage{lmodern}
\usepackage[protrusion=true, expansion=false]{microtype}
\usepackage{csquotes}
\usepackage{textcomp}

\frenchsetup{ThinSpaceInFrenchNumbers=true}

% -----------------------------------------------------------------------------
% Mise en page
% -----------------------------------------------------------------------------
\usepackage[
  top=2.6cm,
  bottom=2.6cm,
  left=2.9cm,
  right=2.9cm,
  headheight=26pt,
  headsep=14pt,
  footskip=28pt
]{geometry}

\usepackage{setspace}
\onehalfspacing

\usepackage{emptypage}
\usepackage{ragged2e}

% -----------------------------------------------------------------------------
% Graphiques, tableaux, listes
% -----------------------------------------------------------------------------
\usepackage{graphicx}
\usepackage{float}
\usepackage{caption}
\usepackage{subcaption}
\usepackage{booktabs}
\usepackage{longtable}
\usepackage{tabularx}
\usepackage{array}
\usepackage{multirow}
\usepackage{makecell}
\usepackage{colortbl}
\usepackage{enumitem}
\usepackage{amsmath, amssymb}
\usepackage{pdflscape}

% Ornements et palette (définit btpblue, btpsky, btplight, btppale, ...)
% =============================================================================
% decorations.tex — Ornements du mémoire (palette bleu clair)
%
% Style : ornements gravés « vectorian » (pgfornament), déclinés en bleu clair.
% Inspiré des mémoires classiques à cadre ornemental, avec un vocabulaire
% graphique propre : cadre à double filet + coins fleuronnés + bandeau répété,
% filets à fleuron pour les chapitres et les sections.
%
% Toutes les commandes publiques sont préfixées \btp...
% =============================================================================

\usepackage{tikz}
\usetikzlibrary{calc}
\usepackage{pgfornament}
\usepackage{eso-pic}
\usepackage{lettrine}

% -----------------------------------------------------------------------------
% Palette — couleur principale #3B5D9C
% Les teintes d'ornement (ciel, clair, pâle, profond) sont mélangées
% à partir de btpblue, pour que cadres, filets et fleurons restent
% dans la même famille.
% -----------------------------------------------------------------------------
\definecolor{btpblue}{HTML}{3B5D9C}           % bleu principal
\colorlet{btpsky}{btpblue!58!white}           % accent (filet intérieur, volutes)
\colorlet{btplight}{btpblue!30!white}          % filets, pieds de page
\colorlet{btppale}{btpblue!10!white}           % fonds d'encadrés et tableaux
\colorlet{btpdeep}{btpblue!58!black}           % titres, numéros de chapitre
\definecolor{btpslate}{RGB}{70, 82, 108}      % gris-bleu (texte secondaire)
\definecolor{btpink}{RGB}{26, 34, 43}          % encre du corps de texte
\definecolor{btpgold}{RGB}{168, 138, 82}      % rehaut discret, accordé au royal

% Alias de compatibilité
\colorlet{maindeco}{btpblue}
\colorlet{decoSlate}{btpslate}
\colorlet{btpaccent}{btpsky}

% -----------------------------------------------------------------------------
% Numéros des ornements vectorian retenus
% (redéfinissables pour changer tout le vocabulaire graphique d'un coup)
% -----------------------------------------------------------------------------
\newcommand{\btpOrnCorner}{63}   % coin fleuronné (page de garde)
\newcommand{\btpOrnCornerLight}{41} % coin à volute (pages courantes)
\newcommand{\btpOrnBand}{88}     % filet fin (bandeaux de cadre)
\newcommand{\btpOrnRule}{85}     % filet à fleuron central (séparateurs)
\newcommand{\btpOrnFlourish}{46} % volute horizontale (chapitres)
\newcommand{\btpOrnSmall}{80}    % petit motif (en-têtes, pieds de page)
\newcommand{\btpOrnLeaf}{75}     % fleuron feuillu (titres liminaires)

% -----------------------------------------------------------------------------
% Séparateurs horizontaux
% -----------------------------------------------------------------------------
% \btpRule[largeur][couleur] : séparateur centré à fleuron
\newcommand{\btpRule}[2][0.55\linewidth]{%
  \begingroup
  \centering
  \pgfornament[width=#1, color=#2]{\btpOrnRule}%
  \par
  \endgroup
}

% Séparateur de section : volute fine sous le titre
\newcommand{\btpSectionRule}{%
  \vspace{-0.35\baselineskip}%
  \noindent
  \textcolor{btplight}{\rule{\linewidth}{0.9pt}}%
  \par\vspace{-0.55\baselineskip}%
  \noindent
  \textcolor{btpblue}{\rule{0.22\linewidth}{2.2pt}}%
  \par\vspace{0.35\baselineskip}%
}

% Fleuron de fin de partie
\newcommand{\btpEndMark}{%
  \par\vspace{0.8em}
  \begin{center}
    \pgfornament[width=3.4cm, color=btpsky]{\btpOrnSmall}
  \end{center}
  \vspace{0.4em}
}

% Fleuron large pour les pages liminaires
\newcommand{\btpLiminaryRule}{%
  \begin{center}
    \pgfornament[width=4.6cm, color=btpblue]{\btpOrnFlourish}
  \end{center}
}

% -----------------------------------------------------------------------------
% Cadre ornemental de la page de garde
% -----------------------------------------------------------------------------
% Double filet + bandeau d'ornements répétés + coins fleuronnés.
% \btpCoverFrame  : à appeler dans l'environnement titlepage.
\newcommand{\btpCoverFrame}{%
  \AddToShipoutPictureFG*{%
    \begin{tikzpicture}[remember picture, overlay]
      % --- filets ---
      \draw[btpblue, line width=1.5pt]
        ($(current page.south west)+(1.15cm,1.15cm)$)
        rectangle
        ($(current page.north east)+(-1.15cm,-1.15cm)$);
      \draw[btpblue!70!white, line width=0.55pt]
        ($(current page.south west)+(1.42cm,1.42cm)$)
        rectangle
        ($(current page.north east)+(-1.42cm,-1.42cm)$);
      % --- bandeaux d'ornements répétés (haut / bas) ---
      \pgfornamentline[color=btpblue!62!white]
        {$(current page.north west)+(2.6cm,-1.285cm)$}
        {$(current page.north east)+(-2.6cm,-1.285cm)$}{9}{\btpOrnBand}
      \pgfornamentline[color=btpblue!62!white]
        {$(current page.south west)+(2.6cm,1.285cm)$}
        {$(current page.south east)+(-2.6cm,1.285cm)$}{9}{\btpOrnBand}
      % --- coins fleuronnés ---
      \node[anchor=north west, inner sep=0] at ($(current page.north west)+(1.18cm,-1.18cm)$)
        {\pgfornament[width=2.35cm, color=btpblue]{\btpOrnCorner}};
      \node[anchor=north east, inner sep=0] at ($(current page.north east)+(-1.18cm,-1.18cm)$)
        {\pgfornament[width=2.35cm, color=btpblue, symmetry=v]{\btpOrnCorner}};
      \node[anchor=south west, inner sep=0] at ($(current page.south west)+(1.18cm,1.18cm)$)
        {\pgfornament[width=2.35cm, color=btpblue, symmetry=h]{\btpOrnCorner}};
      \node[anchor=south east, inner sep=0] at ($(current page.south east)+(-1.18cm,1.18cm)$)
        {\pgfornament[width=2.35cm, color=btpblue, symmetry=c]{\btpOrnCorner}};
    \end{tikzpicture}%
  }%
}

% Cadre allégé pour la page de titre intérieure et les pages liminaires fortes
\newcommand{\btpLightFrame}{%
  \AddToShipoutPictureFG*{%
    \begin{tikzpicture}[remember picture, overlay]
      \draw[btpblue!38!white, line width=0.9pt]
        ($(current page.south west)+(1.3cm,1.3cm)$)
        rectangle
        ($(current page.north east)+(-1.3cm,-1.3cm)$);
      \node[anchor=north west, inner sep=0] at ($(current page.north west)+(1.33cm,-1.33cm)$)
        {\pgfornament[width=1.5cm, color=btpblue]{\btpOrnCornerLight}};
      \node[anchor=south east, inner sep=0] at ($(current page.south east)+(-1.33cm,1.33cm)$)
        {\pgfornament[width=1.5cm, color=btpblue, symmetry=c]{\btpOrnCornerLight}};
    \end{tikzpicture}%
  }%
}

% -----------------------------------------------------------------------------
% Décor optionnel des pages courantes : filet dans la marge extérieure
% Le décor principal des pages courantes est porté par l'en-tête et le pied
% de page (voir preamble.tex) ; ce filet est un supplément facultatif.
% -----------------------------------------------------------------------------
\newcommand{\btpEnableMarginRule}{%
  \AddToShipoutPictureBG{%
    \begin{tikzpicture}[remember picture, overlay]
      \draw[btplight, line width=2.2pt]
        ($(current page.north east)+(-1.15cm,-2.6cm)$) --
        ($(current page.south east)+(-1.15cm,2.6cm)$);
    \end{tikzpicture}%
  }%
}

% Conservé pour compatibilité : n'active plus le filet latéral.
\newcommand{\btpEnablePageDecor}{}

% -----------------------------------------------------------------------------
% Ouverture de chapitre
% -----------------------------------------------------------------------------
% Bloc bleu clair portant le numéro, mention verticale « CHAPITRE »,
% titre en petites capitales, filet à volute.
\newcommand{\btpChapterLabel}{Chapitre}

\newcommand{\btpChapterNumberBlock}[1]{%
  \begin{tikzpicture}[baseline=(blk.center)]
    \node[
      draw=btpblue, line width=1pt,
      fill=btppale,
      rounded corners=1.5pt,
      minimum width=2.05cm, minimum height=2.05cm,
      inner sep=0pt
    ] (blk) {\fontsize{46}{46}\selectfont\bfseries\color{btpdeep}#1};
    \node[
      rotate=90, anchor=south,
      font=\scriptsize\bfseries\scshape,
      text=btpblue,
      inner sep=3pt
    ] at (blk.west) {\btpChapterLabel};
  \end{tikzpicture}%
}

% -----------------------------------------------------------------------------
% Paragraphe d'accroche de chapitre (lettrine)
% -----------------------------------------------------------------------------
\setlength{\DefaultFindent}{2pt}
\renewcommand{\LettrineFontHook}{\color{btpblue}}

% \btpIntro{D}{ans ce chapitre, ...}
%   #1 : lettre initiale (mise en lettrine)
%   #2 : suite du paragraphe d'accroche
\newcommand{\btpIntro}[2]{%
  \begingroup
  \itshape
  \lettrine[lines=3, lhang=0.14, loversize=0.05, findent=3pt, nindent=0pt]%
    {#1}{}#2%
  \par
  \endgroup
  \vspace{0.6em}
  \noindent
  \hbox to \linewidth{\hfil\pgfornament[width=3cm, color=btplight]{\btpOrnSmall}\hfil}%
  \par\vspace{0.5em}
}

% -----------------------------------------------------------------------------
% Bandeau d'image de chapitre (optionnel)
% -----------------------------------------------------------------------------
\newcommand{\btpChapterBanner}[1]{%
  \noindent
  \begin{tikzpicture}
    \clip[rounded corners=3pt]
      (0,0) rectangle (\linewidth, 2.6cm);
    \node[anchor=south west, inner sep=0] at (0,0)
      {\includegraphics[width=\linewidth]{#1}};
    \fill[btpblue, opacity=0.12] (0,0) rectangle (\linewidth, 2.6cm);
  \end{tikzpicture}%
  \par\vspace{0.4cm}%
}

% Alias documenté dans le README : placer juste avant \chapter{...}
\newcommand{\setchapterimage}[1]{\btpChapterBanner{#1}}


\usepackage[
  colorlinks=true,
  linkcolor=btpdeep,
  citecolor=btpblue,
  urlcolor=btpsky,
  filecolor=btpblue,
  bookmarksnumbered=true,
  bookmarksopen=true
]{hyperref}

\usepackage{titlesec}
\usepackage{tocloft}
\usepackage{fancyhdr}
\usepackage[most]{tcolorbox}
\usepackage{listings}
\usepackage{listingsutf8}

% -----------------------------------------------------------------------------
% Chemins d'images
% -----------------------------------------------------------------------------
\graphicspath{%
  {images/page_garde/}%
  {images/liminaire/}%
  {images/chapitre1/}%
  {images/chapitre2/}%
  {images/chapitre3/}%
  {figures/}%
}

\newcommand{\enspylogofile}{logoenspy.png}

% -----------------------------------------------------------------------------
% Colonnes de tableau
% -----------------------------------------------------------------------------
\newcolumntype{L}[1]{>{\raggedright\arraybackslash}p{#1}}
\newcolumntype{C}[1]{>{\centering\arraybackslash}p{#1}}
\newcolumntype{R}[1]{>{\raggedleft\arraybackslash}p{#1}}
\newcolumntype{Y}{>{\raggedright\arraybackslash}X}

\renewcommand{\arraystretch}{1.28}
\setlength{\tabcolsep}{5pt}
\arrayrulecolor{btplight}

% En-tête de tableau sur fond bleu clair
\newcommand{\btpth}[1]{\textcolor{btpdeep}{\bfseries #1}}
\newcommand{\btpTableHeadRule}{\rowcolor{btppale}}

% Environnement de tableau normalisé
\newenvironment{btptable}[2][htbp]
  {\begin{table}[#1]\centering\small\caption{#2}}
  {\end{table}}

% -----------------------------------------------------------------------------
% Légendes
% -----------------------------------------------------------------------------
\captionsetup{
  font=small,
  labelfont={bf, color=btpdeep},
  textfont=it,
  format=plain,
  justification=centering,
  singlelinecheck=false,
  skip=6pt
}
\captionsetup[table]{position=above, skip=8pt}
\captionsetup[figure]{position=below}

% -----------------------------------------------------------------------------
% Titres : chapitre, section, sous-section
% -----------------------------------------------------------------------------
% Blocs de fin de titre construits en primitives TeX : les environnements
% « center » et \par sont interdits dans le code « after » de titlesec.
\newcommand{\btpChapterAfterRule}{%
  \vskip 5pt
  \hbox to \textwidth{\color{btpblue}\leaders\hrule height 1.15pt\hfill}%
  \vskip -7pt
  \hbox to \textwidth{\hfil\pgfornament[width=3.4cm, color=btpsky]{\btpOrnFlourish}}%
  \vskip -6pt
}

\newcommand{\btpChapterAfterFlourish}{%
  \vskip 7pt
  \hbox to \textwidth{\hfil\pgfornament[width=5cm, color=btpblue]{\btpOrnFlourish}\hfil}%
  \vskip -8pt
}

\titleformat{\chapter}[display]
  {\normalfont\filright}
  {\hfill\btpChapterNumberBlock{\thechapter}}
  {6pt}
  {\filright\Large\bfseries\color{btpdeep}}
  [\btpChapterAfterRule]
\titlespacing*{\chapter}{0pt}{-26pt}{24pt}

% Chapitres non numérotés (liminaires, conclusion, annexes)
\titleformat{name=\chapter,numberless}[display]
  {\normalfont\filcenter}
  {}
  {0pt}
  {\filcenter\Large\bfseries\color{btpdeep}}
  [\btpChapterAfterFlourish]
\titlespacing*{name=\chapter,numberless}{0pt}{-16pt}{24pt}

\titleformat{\section}
  {\normalfont\Large\bfseries\color{btpdeep}}
  {\thesection}{0.7em}{}
  [\btpSectionRule]
\titlespacing*{\section}{0pt}{1.5\baselineskip}{0.5\baselineskip}

\titleformat{\subsection}
  {\normalfont\large\bfseries\color{btpblue}}
  {\thesubsection}{0.6em}{}
\titlespacing*{\subsection}{0pt}{1.1\baselineskip}{0.35\baselineskip}

\titleformat{\subsubsection}
  {\normalfont\normalsize\bfseries\itshape\color{btpslate}}
  {\thesubsubsection}{0.5em}{}
\titlespacing*{\subsubsection}{0pt}{0.9\baselineskip}{0.25\baselineskip}

\setcounter{tocdepth}{2}
\setcounter{secnumdepth}{3}

% -----------------------------------------------------------------------------
% Table des matières / figures / tableaux
% -----------------------------------------------------------------------------
\renewcommand{\cftchapfont}{\bfseries\color{btpdeep}}
\renewcommand{\cftchappagefont}{\bfseries\color{btpdeep}}
\renewcommand{\cftsecfont}{\color{btpink}}
\renewcommand{\cftsubsecfont}{\small\color{btpslate}}
\renewcommand{\cftchapleader}{\color{btplight}\cftdotfill{\cftdotsep}}
\renewcommand{\cftsecleader}{\color{btplight}\cftdotfill{\cftdotsep}}
\renewcommand{\cftfigfont}{\color{btpink}}
\renewcommand{\cfttabfont}{\color{btpink}}
\setlength{\cftbeforechapskip}{8pt}

% Titres des tables : même traitement que les chapitres non numérotés
\newcommand{\btpTocTitleFont}{\hspace*{\fill}\Large\bfseries\color{btpdeep}}
\newcommand{\btpTocAfterTitle}{%
  \hspace*{\fill}\par
  \vskip 7pt
  \hbox to \textwidth{\hfil\pgfornament[width=5cm, color=btpblue]{\btpOrnFlourish}\hfil}%
  \vskip 12pt
}
\renewcommand{\cfttoctitlefont}{\btpTocTitleFont}
\renewcommand{\cftloftitlefont}{\btpTocTitleFont}
\renewcommand{\cftlottitlefont}{\btpTocTitleFont}
\renewcommand{\cftaftertoctitle}{\btpTocAfterTitle}
\renewcommand{\cftafterloftitle}{\btpTocAfterTitle}
\renewcommand{\cftafterlottitle}{\btpTocAfterTitle}
\setlength{\cftbeforeloftitleskip}{0pt}
\setlength{\cftbeforelottitleskip}{0pt}
\setlength{\cftbeforetoctitleskip}{0pt}

% -----------------------------------------------------------------------------
% En-têtes et pieds de page
% -----------------------------------------------------------------------------
\pagestyle{fancy}
\fancyhf{}
\fancyhead[L]{\small\itshape\textcolor{btpslate}{\nouppercase{\leftmark}}}
\fancyhead[R]{\small\itshape\textcolor{btpslate}{\nouppercase{\rightmark}}}
\fancyfoot[C]{%
  \pgfornament[width=2.2cm, color=btplight]{\btpOrnSmall}%
  \hspace{0.4em}\small\bfseries\textcolor{btpblue}{\thepage}\hspace{0.4em}%
  \pgfornament[width=2.2cm, color=btplight, symmetry=v]{\btpOrnSmall}%
}
\renewcommand{\headrulewidth}{0.4pt}
\renewcommand{\footrulewidth}{0pt}
\renewcommand{\headrule}{%
  \hbox to\headwidth{\color{btplight}\leaders\hrule height \headrulewidth\hfill}%
}

\fancypagestyle{plain}{%
  \fancyhf{}
  \fancyfoot[C]{%
    \pgfornament[width=2.2cm, color=btplight]{\btpOrnSmall}%
    \hspace{0.4em}\small\bfseries\textcolor{btpblue}{\thepage}\hspace{0.4em}%
    \pgfornament[width=2.2cm, color=btplight, symmetry=v]{\btpOrnSmall}%
  }
  \renewcommand{\headrulewidth}{0pt}
}

\fancypagestyle{empty}{\fancyhf{}\renewcommand{\headrulewidth}{0pt}}

% -----------------------------------------------------------------------------
% Encadrés thématiques
% -----------------------------------------------------------------------------
\tcbuselibrary{breakable, skins}

\tcbset{
  btpbase/.style={
    breakable,
    arc=2.2pt,
    boxrule=0.9pt,
    left=4mm, right=4mm, top=2.4mm, bottom=2.4mm,
    fonttitle=\bfseries,
    enhanced,
    attach boxed title to top left={xshift=6mm, yshift=-2.6mm},
    boxed title style={arc=1.5pt, boxrule=0pt},
  }
}

% Point clé / synthèse
\newtcolorbox{btpbox}[1][Point clé]{
  btpbase,
  colback=btppale, colframe=btpblue,
  coltitle=white,
  boxed title style={colback=btpblue},
  title={#1}
}

% Définition
\newtcolorbox{btpdef}[1][Définition]{
  btpbase,
  colback=white, colframe=btpsky,
  coltitle=white,
  boxed title style={colback=btpsky},
  title={#1}
}

% À retenir / remarque
\newtcolorbox{btpnote}[1][À retenir]{
  btpbase,
  colback=btplight!22, colframe=btpsky,
  coltitle=btpdeep,
  boxed title style={colback=btplight},
  title={#1}
}

% Limite / point de vigilance
\newtcolorbox{btpwarn}[1][Point de vigilance]{
  btpbase,
  colback=btpgold!8, colframe=btpgold,
  coltitle=white,
  boxed title style={colback=btpgold},
  title={#1}
}

% Citation encadrée (problématique, hypothèses)
\newtcolorbox{btpquote}{
  breakable, enhanced,
  colback=btppale, colframe=btpblue,
  boxrule=0pt, leftrule=3.5pt,
  arc=0pt, outer arc=0pt,
  left=5mm, right=4mm, top=3mm, bottom=3mm,
  fontupper=\itshape
}

% -----------------------------------------------------------------------------
% Blocs de code / schémas ASCII
% -----------------------------------------------------------------------------
\lstset{inputencoding=utf8, extendedchars=true}

\lstdefinestyle{btpascii}{
  basicstyle=\ttfamily\scriptsize\color{btpink},
  backgroundcolor=\color{btppale},
  frame=single,
  framerule=0.7pt,
  rulecolor=\color{btplight},
  breaklines=false,
  showstringspaces=false,
  tabsize=2,
  xleftmargin=0.2cm,
  xrightmargin=0.2cm,
  aboveskip=0.8em,
  belowskip=0.8em,
  literate=
    {é}{{\'e}}1 {è}{{\`e}}1 {ê}{{\^e}}1 {ë}{{\"e}}1
    {à}{{\`a}}1 {â}{{\^a}}1 {ç}{{\c c}}1
    {î}{{\^i}}1 {ï}{{\"i}}1 {ô}{{\^o}}1
    {ù}{{\`u}}1 {û}{{\^u}}1 {ü}{{\"u}}1
    {É}{{\'E}}1 {È}{{\`E}}1 {Ê}{{\^E}}1 {À}{{\`A}}1 {Ç}{{\c C}}1
    {°}{{\textdegree}}1 {€}{{\texteuro}}1
}

\lstdefinestyle{btpjson}{
  language=,
  basicstyle=\ttfamily\scriptsize\color{btpink},
  backgroundcolor=\color{btppale},
  keywordstyle=\color{btpdeep}\bfseries,
  stringstyle=\color{btpblue},
  commentstyle=\color{btpslate}\itshape,
  frame=single,
  framerule=0.7pt,
  rulecolor=\color{btplight},
  breaklines=true,
  showstringspaces=false,
  tabsize=2,
  xleftmargin=0.3cm,
  xrightmargin=0.3cm,
  morestring=[b]",
  morekeywords={true, false, null}
}

\lstset{style=btpascii}

% Environnement pour les schémas en art ASCII, avec légende
\newcommand{\btpschema}[2]{%
  \begin{figure}[htbp]
    \centering
    \begin{minipage}{\linewidth}
      \lstinputlisting[style=btpascii]{#1}
    \end{minipage}
    \caption{#2}
  \end{figure}%
}

% -----------------------------------------------------------------------------
% Listes
% -----------------------------------------------------------------------------
\setlist[itemize,1]{label=\textcolor{btpblue}{\small$\blacktriangleright$}, leftmargin=*, itemsep=2pt, topsep=4pt}
\setlist[itemize,2]{label=\textcolor{btpsky}{\small$\bullet$}, itemsep=1pt, topsep=2pt}
\setlist[enumerate,1]{label=\textcolor{btpblue}{\bfseries\arabic*.}, leftmargin=*, itemsep=2pt, topsep=4pt}
\setlist[description]{font=\bfseries\color{btpdeep}, leftmargin=*, itemsep=2pt, topsep=4pt}

% -----------------------------------------------------------------------------
% Styles TikZ partagés par les figures
% -----------------------------------------------------------------------------
\usetikzlibrary{arrows.meta, positioning, calc, fit, shapes.geometric,
                shapes.multipart, backgrounds, chains, decorations.pathreplacing,
                babel}

\tikzset{
  btpbox/.style={
    draw=btpblue, fill=btppale, line width=0.8pt,
    rounded corners=2.5pt, align=center,
    font=\footnotesize, inner sep=4pt,
    minimum height=8mm
  },
  btpboxalt/.style={
    btpbox, fill=btplight!45, draw=btpsky
  },
  btpboxstrong/.style={
    btpbox, fill=btpblue, draw=btpdeep, text=white, font=\footnotesize\bfseries
  },
  btpboxplain/.style={
    btpbox, fill=white, draw=btplight
  },
  btpstep/.style={
    circle, draw=btpblue, fill=white, line width=0.8pt,
    inner sep=1.2pt, minimum size=5.5mm,
    font=\scriptsize\bfseries, text=btpdeep
  },
  btparrow/.style={
    -{Stealth[length=2.2mm, width=1.6mm]},
    draw=btpblue, line width=0.8pt
  },
  btparrowlight/.style={
    -{Stealth[length=2mm, width=1.4mm]},
    draw=btpsky, line width=0.7pt
  },
  btpdashed/.style={
    btparrow, dashed
  },
  btpgroup/.style={
    draw=btpsky, dashed, line width=0.7pt,
    rounded corners=4pt, fill=btplight!14,
    inner sep=6pt
  },
  btplabel/.style={
    font=\scriptsize\itshape, text=btpslate, align=center
  },
  btptitleband/.style={
    fill=btpblue, text=white, font=\small\bfseries\scshape,
    minimum height=7mm, align=center, rounded corners=2pt
  },
  btpdb/.style={
    cylinder, shape border rotate=90, aspect=0.22,
    draw=btpblue, fill=btppale, line width=0.8pt,
    align=center, font=\scriptsize, inner sep=3pt,
    minimum height=11mm, minimum width=15mm
  },
  btpdecision/.style={
    diamond, aspect=1.9, draw=btpblue, fill=btplight!45,
    line width=0.8pt, align=center, font=\scriptsize, inner sep=1pt
  }
}

% Macro homonyme du style TikZ : légende secondaire dans le contenu d'un nœud
\newcommand{\btplabeltext}[1]{{\scriptsize\itshape\color{btpslate}#1}}
\let\btplabel\btplabeltext

% -----------------------------------------------------------------------------
% Chapitres liminaires / annexes non numérotés
% -----------------------------------------------------------------------------
% \btpLiminary{Titre} : chapitre non numéroté, inscrit dans la table des
% matières et repris dans les en-têtes courants.
\newcommand{\btpLiminary}[1]{%
  \chapter*{#1}%
  \addcontentsline{toc}{chapter}{#1}%
  \markboth{#1}{#1}%
}

% -----------------------------------------------------------------------------
% Raccourcis sémantiques
% -----------------------------------------------------------------------------
\newcommand{\sigle}[1]{\textsc{#1}}
\newcommand{\anglais}[1]{\textit{#1}}
\newcommand{\btpterm}[1]{\textbf{#1}}
\newcommand{\oui}{\textcolor{btpblue}{\textbf{Oui}}}
\newcommand{\non}{\textcolor{btpslate}{Non}}
\newcommand{\cotefort}{\textcolor{btpblue}{$\bullet\bullet\bullet$}}
\newcommand{\cotemoyen}{\textcolor{btpsky}{$\bullet\bullet\circ$}}
\newcommand{\cotefaible}{\textcolor{btpslate}{$\bullet\circ\circ$}}

% Étoiles de priorité (remplace les émojis du brouillon)
\newcommand{\prio}[1]{%
  \textcolor{btpblue}{%
    \ifcase#1\relax
      \or $\star$%
      \or $\star\star$%
      \or $\star\star\star$%
      \or $\star\star\star\star$%
      \or $\star\star\star\star\star$%
    \fi
  }%
}

% Marqueurs de maturité / pertinence
\newcommand{\btpMature}{\textcolor{btpblue}{\textbf{Mature}}}
\newcommand{\btpEmergent}{\textcolor{btpgold}{\textbf{Émergent}}}
\newcommand{\btpEnDev}{\textcolor{btpsky}{\textbf{En développement}}}

% Numérotation des figures et tableaux par chapitre
\renewcommand{\thefigure}{\thechapter.\arabic{figure}}
\renewcommand{\thetable}{\thechapter.\arabic{table}}
